% Add \usepackage{booktabs,adjustbox} to your manuscript preamble.
\begin{table*}[t]
\centering
\small
\caption{BCN20000 study partitions}
\label{tab:01-dataset-partitions}
\begin{adjustbox}{max width=\textwidth}
\begin{tabular}{lll}
\toprule
\textbf{Partition} & \textbf{Images} & \textbf{Leakage groups} \\
\midrule
Development & 7,813 & 2,286 \\
Original validation (excluded) & 1,973 & 571 \\
Reserved test & 2,578 & 715 \\
Total & 12,364 & 3,572 \\
\bottomrule
\end{tabular}
\end{adjustbox}
\par\smallskip
\begin{minipage}{\textwidth}\footnotesize Partition counts from Cell 111 V2. Development used five group-aware folds. Reserved results are from the structurally audited same-dataset snapshot; historical split integrity and historical nonaccess were not established. No external validation or test-based selection.\end{minipage}
\end{table*}

\begin{table*}[t]
\centering
\small
\caption{Reserved-test class distribution}
\label{tab:02-reserved-class-support}
\begin{adjustbox}{max width=\textwidth}
\begin{tabular}{llll}
\toprule
\textbf{Class} & \textbf{Images} & \textbf{Groups} & \textbf{Image share (\%)} \\
\midrule
AK & 129 & 48 & 5.00 \\
BCC & 580 & 196 & 22.50 \\
BKL & 241 & 70 & 9.35 \\
DF & 29 & 8 & 1.12 \\
MEL & 632 & 105 & 24.52 \\
NV & 842 & 255 & 32.66 \\
SCC & 94 & 25 & 3.65 \\
VASC & 31 & 8 & 1.20 \\
\bottomrule
\end{tabular}
\end{adjustbox}
\par\smallskip
\begin{minipage}{\textwidth}\footnotesize Image counts and group counts are different units; the reserved test has 2,578 images and 715 leakage groups.\end{minipage}
\end{table*}

\begin{table*}[t]
\centering
\small
\caption{Controlled architectural comparison}
\label{tab:03-ablation-design}
\begin{adjustbox}{max width=\textwidth}
\begin{tabular}{llllll}
\toprule
\textbf{Arm} & \textbf{Architecture} & \textbf{Clinical inputs} & \textbf{GATv2} & \textbf{Adaptive mechanism} & \textbf{Head parameters} \\
\midrule
H0 & Image-only & No & No & No & 126,856 \\
H1 & Non-graph multimodal & Yes & No & No & 195,665 \\
H2 & Non-adaptive GATv2 & Yes & Yes & No & 195,665 \\
H3 & Adaptive GATv2 (proposed) & Yes & Yes & Yes & 195,665 \\
H4 & Adaptive non-graph control & Yes & No & Yes & 195,665 \\
\bottomrule
\end{tabular}
\end{adjustbox}
\par\smallskip
\begin{minipage}{\textwidth}\footnotesize All arms use the frozen 1,440-dimensional Swin-T feature representation and the same four-epoch final training budget. Counts are trainable head parameters, not full-backbone counts.\end{minipage}
\end{table*}

\begin{table*}[t]
\centering
\small
\caption{Grouped five-fold out-of-fold development performance}
\label{tab:04-development-oof}
\begin{adjustbox}{max width=\textwidth}
\begin{tabular}{llllll}
\toprule
\textbf{Arm} & \textbf{Accuracy (\%)} & \textbf{Macro-F1 (\%)} & \textbf{Balanced acc. (\%)} & \textbf{Fold Macro-F1 mean (\%)} & \textbf{Fold SD (pp)} \\
\midrule
H0 & 55.80 & 37.38 & 37.93 & 36.81 & 4.05 \\
H1 & 56.87 & 37.34 & 37.15 & 37.34 & 4.47 \\
H2 & 57.12 & 38.41 & 38.75 & 37.87 & 3.77 \\
H3 & 57.03 & 37.73 & 38.24 & 37.57 & 3.24 \\
H4 & 56.87 & 36.37 & 36.63 & 36.13 & 2.82 \\
\bottomrule
\end{tabular}
\end{adjustbox}
\par\smallskip
\begin{minipage}{\textwidth}\footnotesize Pooled OOF metrics from 7,813 development images in 2,286 groups. Fold mean and sample SD are separate from pooled Macro-F1. Development results are not reserved-test results.\end{minipage}
\end{table*}

\begin{table*}[t]
\centering
\small
\caption{Reserved-test performance under natural metadata availability}
\label{tab:05-reserved-natural-performance}
\begin{adjustbox}{max width=\textwidth}
\begin{tabular}{lllllll}
\toprule
\textbf{Arm} & \textbf{Accuracy (\%)} & \textbf{Macro-F1 (\%)} & \textbf{Balanced acc. (\%)} & \textbf{Macro ROC-AUC} & \textbf{MEL recall (\%)} & \textbf{SCC recall (\%)} \\
\midrule
H0 & 58.81 & 40.59 & 41.98 & 0.8487 & 62.50 & 36.17 \\
H1 & 60.43 & 40.30 & 41.02 & 0.8568 & 64.87 & 36.17 \\
H2 & 61.09 & 40.07 & 40.81 & 0.8566 & 65.35 & 38.30 \\
H3 & 61.79 & 41.65 & 43.11 & 0.8593 & 68.99 & 46.81 \\
H4 & 60.98 & 41.71 & 43.23 & 0.8579 & 68.99 & 42.55 \\
\bottomrule
\end{tabular}
\end{adjustbox}
\par\smallskip
\begin{minipage}{\textwidth}\footnotesize Macro ROC-AUC is one-vs-rest on a 0--1 scale. H3 is the proposed arm, not a test-selected winner. Reserved results are from the structurally audited same-dataset snapshot; historical split integrity and historical nonaccess were not established. No external validation or test-based selection.\end{minipage}
\end{table*}

\begin{table*}[t]
\centering
\small
\caption{Natural-condition performance with paired group-bootstrap uncertainty}
\label{tab:06-reserved-uncertainty}
\begin{adjustbox}{max width=\textwidth}
\begin{tabular}{lllll}
\toprule
\textbf{Arm} & \textbf{Accuracy (\%)} & \textbf{95\% CI} & \textbf{Macro-F1 (\%)} & \textbf{95\% CI} \\
\midrule
H0 & 58.81 & [55.50, 62.02] & 40.59 & [36.78, 44.68] \\
H1 & 60.43 & [57.08, 63.67] & 40.30 & [36.39, 44.32] \\
H2 & 61.09 & [57.80, 64.30] & 40.07 & [36.58, 43.64] \\
H3 & 61.79 & [58.47, 65.02] & 41.65 & [37.71, 45.75] \\
H4 & 60.98 & [57.68, 64.19] & 41.71 & [37.61, 45.93] \\
\bottomrule
\end{tabular}
\end{adjustbox}
\par\smallskip
\begin{minipage}{\textwidth}\footnotesize Descriptive 95\% percentile intervals, 10,000 saved class-stratified group-bootstrap draws. These are not the multiplicity-adjusted primary contrast intervals.\end{minipage}
\end{table*}

\begin{table*}[t]
\centering
\small
\caption{Proposed H3: per-class reserved-test performance}
\label{tab:07-h3-per-class}
\begin{adjustbox}{max width=\textwidth}
\begin{tabular}{lllllll}
\toprule
\textbf{Class} & \textbf{Images} & \textbf{Precision (\%)} & \textbf{Recall (\%)} & \textbf{Specificity (\%)} & \textbf{F1 (\%)} & \textbf{ROC-AUC} \\
\midrule
AK & 129 & 27.61 & 28.68 & 96.04 & 28.14 & 0.8591 \\
BCC & 580 & 69.08 & 62.41 & 91.89 & 65.58 & 0.9040 \\
BKL & 241 & 36.79 & 29.46 & 94.78 & 32.72 & 0.7764 \\
DF & 29 & 11.36 & 17.24 & 98.47 & 13.70 & 0.7893 \\
MEL & 632 & 70.66 & 68.99 & 90.70 & 69.82 & 0.8822 \\
NV & 842 & 72.59 & 75.18 & 86.23 & 73.86 & 0.9019 \\
SCC & 94 & 29.33 & 46.81 & 95.73 & 36.07 & 0.8814 \\
VASC & 31 & 11.36 & 16.13 & 98.47 & 13.33 & 0.8797 \\
\bottomrule
\end{tabular}
\end{adjustbox}
\par\smallskip
\begin{minipage}{\textwidth}\footnotesize Sensitivity equals recall. AUC is one-vs-rest (0--1). Point estimates only; small class supports limit interpretation.\end{minipage}
\end{table*}

\begin{table*}[t]
\centering
\small
\caption{Missing-metadata stress-test performance}
\label{tab:08-metadata-robustness}
\begin{adjustbox}{max width=\textwidth}
\begin{tabular}{lllll}
\toprule
\textbf{Arm} & \textbf{Natural F1 (\%)} & \textbf{Worst F1 (\%)} & \textbf{Worst drop (pp)} & \textbf{All-missing F1 (\%)} \\
\midrule
H1 & 40.30 & 37.77 & 2.52 & 37.83 \\
H2 & 40.07 & 37.32 & 2.76 & 37.32 \\
H3 & 41.65 & 39.24 & 2.41 & 39.24 \\
H4 & 41.71 & 39.62 & 2.09 & 40.29 \\
\bottomrule
\end{tabular}
\end{adjustbox}
\par\smallskip
\begin{minipage}{\textwidth}\footnotesize Worst is the minimum over natural plus seven forced-removal scenarios (M00--M07). Smaller degradation is better; larger absolute F1 is better. H0 is image-only and does not require metadata stress testing.\end{minipage}
\end{table*}

\begin{table*}[t]
\centering
\small
\caption{Predeclared primary contrasts with Bonferroni-adjusted confidence intervals}
\label{tab:09-primary-contrasts}
\begin{adjustbox}{max width=\textwidth}
\begin{tabular}{lllll}
\toprule
\textbf{Contrast} & \textbf{Endpoint} & \textbf{Difference (pp)} & \textbf{98.33\% CI (pp)} & \textbf{Positive lower bound} \\
\midrule
H1-H0 & Natural Macro-F1 & -0.29 & [-3.19, +2.42] & No \\
H3-H2 & Worst-case Macro-F1 & +1.92 & [-0.48, +4.95] & No \\
H3-H4 & Worst-case Macro-F1 & -0.38 & [-1.59, +0.85] & No \\
\bottomrule
\end{tabular}
\end{adjustbox}
\par\smallskip
\begin{minipage}{\textwidth}\footnotesize 10,000 paired class-stratified group-bootstrap draws. Worst-case endpoint recomputed within each draw. None of the three primary superiority criteria was met; this does not establish equivalence.\end{minipage}
\end{table*}

\begin{table*}[t]
\centering
\small
\caption{Exploratory image-occlusion faithfulness on development examples}
\label{tab:10-xai-occlusion}
\begin{adjustbox}{max width=\textwidth}
\begin{tabular}{llll}
\toprule
\textbf{Intervention} & \textbf{Cases} & \textbf{Mean target-logit drop} & \textbf{Top \textgreater{} comparator (cases)} \\
\midrule
Top-ranked 4 tiles & 16 & 1.7371 & N/A \\
Fixed random 4 tiles & 16 & 0.6187 & 15 \\
Bottom-ranked 4 tiles & 16 & -0.2007 & 16 \\
\bottomrule
\end{tabular}
\end{adjustbox}
\par\smallskip
\begin{minipage}{\textwidth}\footnotesize Sixteen class-stratified training examples, one per group. Positive drop means the original predicted-class logit decreased after deletion. This is not independent clinical explanation validation.\end{minipage}
\end{table*}

\begin{table*}[t]
\centering
\small
\caption{Exploratory metadata-removal analysis on development examples}
\label{tab:11-xai-metadata}
\begin{adjustbox}{max width=\textwidth}
\begin{tabular}{lll}
\toprule
\textbf{Removed clinical information} & \textbf{Changed cases} & \textbf{Mean target-logit drop} \\
\midrule
Age & 16 & 0.9721 \\
Sex & 16 & 0.8469 \\
Anatomical site & 16 & 1.0929 \\
All three fields & 16 & 2.4583 \\
\bottomrule
\end{tabular}
\end{adjustbox}
\par\smallskip
\begin{minipage}{\textwidth}\footnotesize Means are over cases whose observation changed, not necessarily all 16 examples. Interventions are model perturbations, not causal clinical evidence.\end{minipage}
\end{table*}

\begin{table*}[t]
\centering
\small
\caption{Clinical metadata availability in the reserved test}
\label{tab:12-metadata-coverage}
\begin{adjustbox}{max width=\textwidth}
\begin{tabular}{llll}
\toprule
\textbf{Field} & \textbf{Observed images} & \textbf{Missing images} & \textbf{Observed (\%)} \\
\midrule
Age & 2,558 & 20 & 99.22 \\
Sex & 2,558 & 20 & 99.22 \\
Anatomical site & 2,551 & 27 & 98.95 \\
\bottomrule
\end{tabular}
\end{adjustbox}
\par\smallskip
\begin{minipage}{\textwidth}\footnotesize Observed counts are field-specific; overlapping missingness is possible. Availability at prospective clinical prediction time was not validated.\end{minipage}
\end{table*}

\begin{table*}[t]
\centering
\small
\caption{Uncalibrated probability performance under natural metadata availability}
\label{tab:13-calibration}
\begin{adjustbox}{max width=\textwidth}
\begin{tabular}{llll}
\toprule
\textbf{Arm} & \textbf{ECE (15 bins)} & \textbf{Multiclass Brier sum} & \textbf{Negative log-likelihood} \\
\midrule
H0 & 0.0773 & 0.5465 & 1.1514 \\
H1 & 0.1161 & 0.5375 & 1.1461 \\
H2 & 0.1113 & 0.5348 & 1.1435 \\
H3 & 0.1180 & 0.5411 & 1.1750 \\
H4 & 0.1244 & 0.5451 & 1.1791 \\
\bottomrule
\end{tabular}
\end{adjustbox}
\par\smallskip
\begin{minipage}{\textwidth}\footnotesize Lower values are better. Point estimates from saved probabilities; no post-hoc calibration or threshold fitting.\end{minipage}
\end{table*}
